\documentclass[11pt]{article}

\usepackage[T1]{fontenc}
\usepackage[utf8]{inputenc}
\usepackage[a4paper,margin=1in]{geometry}
\usepackage{booktabs}
\usepackage{graphicx}
\usepackage{xcolor}
\usepackage{titlesec}
\usepackage{enumitem}
\usepackage[hidelinks]{hyperref}

\definecolor{accent}{HTML}{C0392B}
\definecolor{blu}{HTML}{2C6FBB}
\definecolor{ink}{HTML}{2B2B2B}

\titleformat{\section}{\Large\bfseries\color{ink}}{\thesection}{0.6em}{}
\titleformat{\subsection}{\large\bfseries\color{blu}}{\thesubsection}{0.6em}{}
\setlist{itemsep=2pt,topsep=3pt}
\graphicspath{{figures/}}

\newcommand{\artist}[1]{\textit{#1}}

\begin{document}

\begin{titlepage}
\centering
\vspace*{3cm}
{\color{accent}\rule{\textwidth}{2pt}}\\[0.8em]
{\Huge\bfseries A Portrait of a Predictor}\\[0.4em]
{\Large What your next-track model hears coming}\\[0.6em]
{\large\itshape Un retrato de quien predice --- lo que tu modelo del próximo tema oye venir}\\[0.8em]
{\color{accent}\rule{\textwidth}{2pt}}\\[2cm]
{\large A data portrait derived from 91{,}827 listening steps, the concepts a next-track\\
GRU learned inside its own hidden state, and GPT-labeled sparse-autoencoder atoms.}\\[1cm]
{\normalsize Spotify Next-Track Lab \quad$\cdot$\quad sequence dataset \texttt{seq-20260718-222807}}\\[0.3em]
{\normalsize 2026-07-23}
\vfill
\end{titlepage}

% ============================================================
\section*{How to read this document}
This report is bilingual. \textbf{Part~I is in English}; \textbf{Part~II es la
versión en español} (identical content). It reads one thing: not \emph{what} you
listen to, but what a machine trained to guess your \emph{next} track has learned
about you. It is built from (1)~91{,}827 listening steps (each: the session so
far $\rightarrow$ the track that actually came next); (2)~the champion blend's
recurrent base learner, a GRU (\texttt{best-seq-nexttrack-20260718-230948-d3d4a})
that predicts the next track's latent and ranks the library by it; and (3)~a
sparse-autoencoder (SAE) probe of that GRU's hidden state, whose learned
``atoms'' were named by GPT (\texttt{gpt-4o-mini}). Method and caveats are in the
Provenance appendix --- read it: the model has \emph{genuinely} learned next-track
structure, but it stores it in a \emph{dense, entangled} way, and that honesty is
part of the picture.

\clearpage
\begin{center}{\color{blu}\Large\bfseries PART I \quad$\cdot$\quad ENGLISH}\end{center}
\vspace{0.5em}

\section{The one-paragraph portrait}
Ask this model what you will play next and it thinks, above all, in
\textbf{artists and their sonic neighbourhoods}. Its strongest internal features
are dedicated detectors for the artists you loop hardest --- \artist{DIIV},
\artist{Rammstein}, \artist{Thievery Corporation}, \artist{Aphex Twin},
\artist{Four Tet} --- and for a handful of genre textures: \emph{alternative
metal}, \emph{Brooklyn / indie rock}, \emph{album rock}, \emph{Düsseldorf
electronic}, \emph{ambient}, and a clear seam of \emph{argentine rock}. From the
session so far it can linearly guess the \emph{genre} of your very next track far
above chance. It has learned the same listener the aggregate portrait found ---
electronic core, rock heart, Argentine thread --- but arrived at it sideways, by
learning what tends to \emph{follow} what in a session.

\section{What it predicts}
The model ranks your whole library for the next slot. On held-out sessions the
true next track lands in its top~10 \textbf{15.7\%} of the time
(\texttt{recall@10}), and when it misses the exact track it usually lands
something that \emph{sounds} right: its top~10 scores \textbf{0.48}
on \texttt{music@10} (sonic similarity to the truth). Most tellingly, the next
track's \textbf{genre is linearly decodable from the GRU's hidden state at
AUC~0.79} (41.7\% top-1 across 25~genres, against a 0.7\% chance rate) --- the
recurrent state genuinely carries where the session is heading.

\begin{center}\includegraphics[width=0.62\textwidth]{mind_decodability.pdf}\end{center}

\section{The concepts it formed}
The GRU doesn't store ``predict DIIV.'' It learns abstract feature detectors in
its 256-dimensional hidden state; the SAE pulls those apart into individual
``atoms,'' and GPT named each from the next-tracks that most excite it. The
strongest atoms are crisp and namable:

\begin{center}
\begin{tabular}{@{}rll@{}}
\toprule
sep ($\sigma$) & Concept (GPT label) & Type \\
\midrule
$3.8$ & DIIV --- Brooklyn indie vibe & artist \\
$3.5$ & Rammstein --- live performances & artist \\
$3.4$ & Rammstein --- alternative-metal pattern & artist \\
$3.0$ & Pink Floyd --- album-rock pattern & genre \\
$2.8$ & Thievery Corporation --- downtempo & artist \\
$2.8$ & Indie-rock favourites & genre \\
$2.8$ & DIIV --- indie soundscape & artist \\
$2.7$ & Aphex Twin --- ambient detection & artist \\
$2.7$ & Four Tet --- electronica repetition & artist \\
\bottomrule
\end{tabular}
\end{center}

\begin{center}\includegraphics[width=0.86\textwidth]{mind_concepts.pdf}\end{center}

Each ``sep'' is how many standard deviations an atom's activation rises right
before a next track of that artist/genre, versus every other step --- a clean
$3.8\sigma$ for the DIIV detector is a strong, specific feature.

\section{Who it thinks about most}
The model spends its sharpest features on the artists that dominate your
sessions. \textbf{Rammstein alone accounts for several of the top atoms}; DIIV
gets two. This is the recurrent echo of heavy rotation: the more an artist
anchors your sessions, the more dedicated hidden capacity the GRU devotes to
seeing them coming. Some of these atoms also track \emph{sonic continuity} --- the
Rammstein detectors carry a mood correlation of $\approx +0.2$, firing when the
next track is not just same-artist but same-\emph{sound}, while the DIIV detector
is almost purely identity ($\approx 0.0$).

\section{What it captures --- and what it drops}
Reassuringly, the model leaves almost nothing on the table: \textbf{all 25 of the
frequent next-track genres are represented} by at least one dedicated atom, and
\textbf{zero} frequent next-item concepts are dropped. The genres it separates
most sharply are exactly your pillars:

\begin{center}\includegraphics[width=0.82\textwidth]{mind_genre_repr.pdf}\end{center}

\subsection{The honest caveat: dense, not clean}
The catch is \emph{how} it stores all this. The hidden state is
\textbf{dense, not sparse}: the SAE reconstructs it near-perfectly
($99\%$ of variance), but $\approx 418$ of $512$ atoms fire on a typical step. So
while the strongest atoms are namable, next-track information is \emph{smeared
across nearly the whole layer} rather than isolated into a few tidy features. Read
the concepts above as ``the dominant artists and genres each have a strong
dedicated direction,'' not ``the model is a clean dictionary of your taste.'' This
matches the standing lab finding that SAEs on this corpus don't yield crisp
monosemantic atoms at this data scale --- the labels are suggestive readings, and
the decodability numbers ($\S2$) are the solid ground.

\section{The portrait in one line}
\begin{quote}\itshape
A next-track model that thinks in artists and their sonic neighbourhoods ---
sharpest on the ones you loop (DIIV, Rammstein, Thievery Corporation, Aphex Twin,
Four Tet), fluent in every genre from alternative metal to Düsseldorf electronic
to argentine rock, able to call the next track's genre far above chance --- but
holding it all in a dense, distributed hidden state rather than a tidy shelf of
concepts.
\end{quote}

\clearpage
\begin{center}{\color{blu}\Large\bfseries PARTE II \quad$\cdot$\quad ESPAÑOL}\end{center}
\vspace{0.5em}

\section{El retrato en un párrafo}
Si le preguntas a este modelo qué vas a reproducir a continuación, piensa, ante
todo, en \textbf{artistas y sus vecindarios sonoros}. Sus rasgos internos más
fuertes son detectores dedicados a los artistas que más repites ---
\artist{DIIV}, \artist{Rammstein}, \artist{Thievery Corporation}, \artist{Aphex
Twin}, \artist{Four Tet}--- y a unas pocas texturas de género: \emph{alternative
metal}, \emph{indie} de Brooklyn, \emph{album rock}, electrónica de Düsseldorf,
\emph{ambient} y una veta clara de \emph{rock argentino}. A partir de la sesión
hasta ese punto puede adivinar linealmente el \emph{género} del próximo tema muy
por encima del azar. Ha aprendido al mismo oyente que encontró el retrato
agregado --- núcleo electrónico, corazón de rock, hebra argentina--- pero llegó de
costado, aprendiendo qué suele \emph{seguir} a qué dentro de una sesión.

\section{Lo que predice}
El modelo ordena toda tu biblioteca para el próximo lugar. En sesiones reservadas,
el próximo tema real cae en su top~10 el \textbf{15,7\,\%} de las veces
(\texttt{recall@10}), y cuando falla el tema exacto suele acertar algo que
\emph{suena} bien: su top~10 marca \textbf{0,48} en \texttt{music@10} (similitud
sonora con lo verdadero). Y lo más revelador: el \textbf{género del próximo tema
es linealmente decodificable desde el estado oculto del GRU con AUC~0,79} (41,7\,\%
de acierto top-1 entre 25~géneros, frente a un azar del 0,7\,\%) --- el estado
recurrente sí lleva hacia dónde va la sesión.

\begin{center}\includegraphics[width=0.62\textwidth]{mind_decodability.pdf}\end{center}
\begin{center}\small\itshape (Etiquetas de la figura en inglés: ``accuracy'' $=$ acierto; ``chance'' $=$ azar.)\end{center}

\section{Los conceptos que formó}
El GRU no guarda ``predecir DIIV''. Aprende detectores de rasgos abstractos en su
estado oculto de 256 dimensiones; el SAE los separa en ``átomos'' individuales, y
GPT nombró cada uno a partir de los próximos temas que más lo activan. Los átomos
más fuertes son nítidos y nombrables:

\begin{center}
\begin{tabular}{@{}rll@{}}
\toprule
sep ($\sigma$) & Concepto (etiqueta GPT) & Tipo \\
\midrule
$3{,}8$ & DIIV --- \emph{indie} de Brooklyn & artista \\
$3{,}5$ & Rammstein --- en vivo & artista \\
$3{,}4$ & Rammstein --- patrón \emph{alt-metal} & artista \\
$3{,}0$ & Pink Floyd --- patrón \emph{album rock} & género \\
$2{,}8$ & Thievery Corporation --- \emph{downtempo} & artista \\
$2{,}8$ & Favoritos del \emph{indie rock} & género \\
$2{,}8$ & DIIV --- paisaje \emph{indie} & artista \\
$2{,}7$ & Aphex Twin --- \emph{ambient} & artista \\
$2{,}7$ & Four Tet --- repetición \emph{electronica} & artista \\
\bottomrule
\end{tabular}
\end{center}

\begin{center}\includegraphics[width=0.86\textwidth]{mind_concepts.pdf}\end{center}

Cada ``sep'' es cuántas desviaciones estándar sube la activación de un átomo justo
antes de un próximo tema de ese artista/género, frente al resto de los pasos --- un
$3{,}8\sigma$ limpio para el detector de DIIV es un rasgo fuerte y específico.

\section{En quién piensa más}
El modelo gasta sus rasgos más afilados en los artistas que dominan tus sesiones.
\textbf{Rammstein solo se lleva varios de los átomos principales}; DIIV, dos. Es el
eco recurrente de la reescucha intensa: cuanto más ancla un artista tus sesiones,
más capacidad oculta le dedica el GRU a verlo venir. Algunos de esos átomos también
siguen la \emph{continuidad sonora} --- los detectores de Rammstein llevan una
correlación de humor de $\approx +0{,}2$, activándose cuando el próximo tema no solo
es del mismo artista sino del mismo \emph{sonido}, mientras que el de DIIV es casi
pura identidad ($\approx 0{,}0$).

\section{Lo que captura --- y lo que deja pasar}
Tranquiliza que el modelo casi no deja nada sobre la mesa: \textbf{los 25 géneros
frecuentes del próximo tema están representados} por al menos un átomo dedicado, y
se descartan \textbf{cero} conceptos frecuentes. Los géneros que separa con más
nitidez son justamente tus pilares:

\begin{center}\includegraphics[width=0.82\textwidth]{mind_genre_repr.pdf}\end{center}

\subsection{La advertencia honesta: denso, no limpio}
El detalle es \emph{cómo} guarda todo esto. El estado oculto es \textbf{denso, no
disperso}: el SAE lo reconstruye casi a la perfección ($99\%$ de la varianza), pero
$\approx 418$ de $512$ átomos se activan en un paso típico. Así que, aunque los
átomos más fuertes sean nombrables, la información del próximo tema está
\emph{repartida por casi toda la capa} en vez de aislada en unos pocos rasgos
limpios. Lee los conceptos de arriba como ``los artistas y géneros dominantes
tienen cada uno una dirección dedicada fuerte'', no como ``el modelo es un
diccionario limpio de tu gusto''. Esto coincide con el hallazgo estable del
laboratorio de que los SAE en este corpus no dan átomos monosemánticos nítidos a
esta escala de datos --- las etiquetas son lecturas sugerentes, y los números de
decodificación ($\S2$) son el terreno firme.

\section{El retrato en una línea}
\begin{quote}\itshape
Un modelo del próximo tema que piensa en artistas y sus vecindarios sonoros --- más
afilado en los que repites (DIIV, Rammstein, Thievery Corporation, Aphex Twin,
Four Tet), fluido en cada género del \emph{alternative metal} a la electrónica de
Düsseldorf y el rock argentino, capaz de anticipar el género del próximo tema muy
por encima del azar --- pero guardándolo todo en un estado oculto denso y
distribuido en vez de un estante ordenado de conceptos.
\end{quote}

\clearpage
\section*{Provenance \& honest caveats \quad$\cdot$\quad Procedencia y advertencias}

\textbf{Data.} A single user's Spotify listening history, sessionized into the
sequence dataset \texttt{seq-20260718-222807}. The per-model SAE ran over
\textbf{91{,}827} \emph{(session, step)} rows (73{,}632 train / 18{,}195 test);
each row is the hidden state at one step paired with the track that actually came
next. Aggregate concept statistics are computed directly over these rows.

\textbf{Model.} \texttt{best-seq-nexttrack-20260718-230948-d3d4a}, a
single-layer GRU (256-d hidden) --- the recurrent base learner of the promoted
GRU$\times$Markov blend champion. It encodes the session prefix and ranks the
library by cosine to its predicted next-track latent
(\texttt{recall@10}~$=0.157$, \texttt{music@10}~$=0.48$).

\textbf{Interpretability.} A sparse autoencoder (512 atoms, $2\times$ overcomplete,
$\ell_1=0.0015$, 40 epochs) trained on the GRU's per-step hidden activations;
analysis \texttt{interp-run-20260723-204848-4e89f-msae} (reopen it under the UI's
\emph{Saved analyses} tool). ``Concepts'' are read against the \emph{next item}:
an atom's association is its mean-activation separation (in $\sigma$) for a
next-track artist / genre / album class. Next-item genre decodability is a
logistic probe on the raw activations (top-1 accuracy + macro-AUC vs. a majority
baseline). The 12 surfaced atoms per layer had their top-activating next-tracks
sent to \textbf{GPT (\texttt{gpt-4o-mini})} for a concept label.

\textbf{The honest caveat.} The GRU has genuinely learned next-track structure ---
genre is decodable at AUC~0.79 and every frequent concept is represented --- but the
representation is \emph{dense}: $\approx 418$ of $512$ atoms fire per step and the
code reconstructs $99\%$ of the layer's variance, so ``concepts'' are strong
\emph{directions} in an entangled space, not clean monosemantic units. The SAE is a
discovery lens, not a decision rule; the GPT labels are automated readings of
10~tracks each and occasionally blur. The decodability and coverage numbers are
solid; the atom labels are illustrative.

\vfill
\begin{center}\small
Generated from the Spotify Next-Track Lab
$\cdot$ analysis \texttt{interp-run-20260723-204848-4e89f-msae}
$\cdot$ figures: \texttt{docs/figures/make\_mind\_figures.py}
\end{center}

\end{document}
